
\section{Suspended sand settling into a dredged trench: van Rijn's flume}

Delft Hydraulics flume experiments reported by van Rijn \cite{vanRijn1986b}
(Figs.~16--17): a flow $h_0 = 0.39$~m deep at $\bar u_0 = 0.51$~m/s carrying
160~$\mu$m sand in equilibrium (0.04~kg/s/m in total, 0.03 in suspension and
0.01 as bedload; suspended size 130~$\mu$m, $w_s = 0.013$~m/s; $k_s = 0.025$~m)
crosses a trench cut across the 0.5~m wide flume. The flow slows over the
trench, the suspended sand settles onto the downstream side slope and the
floor, and the trench fills and migrates downstream. The bed profiles
measured after 15 hours for side slopes of 1:3, 1:7 and 1:10 (the ``x''
markers of Fig.~16, digitised here) are the data.

\anuga{} runs the approach flow as normal flow on a plane with the Manning
$n$ of $k_s$, held by Dirichlet boundaries, with the trench cut into the
plane 5~m from the inflow, and the inflow carrying the measured suspended
supply. One sand fraction carries the measured settling velocity
(Ferguson--Church equivalent diameter 138~$\mu$m), the de Leeuw et al.
(2020) entrainment law with the Rouse near-bed ratio referenced at $0.1h$,
Wong--Parker bedload with open inflow and outflow, and an evolving bed. As
van Rijn did for his own model (``the constant of Eq.~26 was adjusted
somewhat to give a suspended load of 0.03~kg/sm at the inlet''), the
entrainment constant is calibrated so that the approach flow is in
equilibrium at the measured supply, by short pre-runs on the plane without
the trench; with the de Leeuw law the uncalibrated equilibrium is already
0.0305~kg/s/m, so the factor is 0.98. What the trench then tests is the
response of the transport to the non-uniform flow: the settling lag, the
deposition on the downstream slope and floor, and the Exner coupling.

The flume is one-dimensional, so the mesh carries two rows of 0.1~m cells
across it, and the 15 hours of bed evolution run under a morphological
acceleration factor of 10: the bed change of every step is multiplied by
ten, the suspension, which adapts over a settling length of about 20~m in
under a minute, is left alone, and the run covers 1.5 hours of flow.
Wong--Parker bedload is scaled to the measured 0.01~kg/s/m at the inlet,
as van Rijn did with his own bedload constant, and the inlet carries that
bedload as a prescribed flux, as the flume was fed at a set rate. The
default zero-gradient import, which equals the inflow cell's own export,
feeds back on itself here: a cell that aggrades under the fixed inflow
stage sees its transport and its import rise, and within a few hours the
reach aggrades without bound. A rigid approach strip cures that but starves
the first erodible cell behind it of bedload, which then scours a 25--30~cm
hole that feeds the trench; the prescribed supply avoids both, and the two
compute modes then agree to 0.1~mm over the 15 hours.

The pass criterion is the bed level after 15 hours at the measured points:
RMS error below 0.03~m (the fill is 0.08--0.10~m deep over 7~m of flume,
and van Rijn's own profile model matched it to about 0.01~m), and the
deepest remaining point within 1.5~m of the measured position. The gentle
1:10 trench, the one a depth-averaged model is expected to get, runs
routinely; the 1:3 and 1:7 trenches, where the flow separates over the lip
and van Rijn's own depth-averaged attempts failed, join it under the long
option. \anuga{} reproduces the shape of the measured profile and the
position of its deepest point but fills the trench about 25\,\% too fast:
the near-bed concentration responds at once to the slower flow over the
trench, where in the flume the sediment high in the water column takes time
to settle through it. The adaptation lag of Galappatti and Vreugdenhil
(\texttt{set\_deposition(adaptation='armanini')}) over-corrects this, so the
measured profile lies between the two, and the carried near-bed ratio
(\texttt{adaptation='carried'}) recovers part of the gap, 3.5, 3.1 and
2.7~cm RMS against 3.9, 3.2 and 2.85 without it; the case runs without
either, so that the closure it validates is the default one.

\begin{figure}
\begin{center}
\includegraphics[width=0.9\textwidth]{bed_profile_test3.png}
\end{center}
\caption{Test 3, side slopes 1:10: bed level at three times against the
measured profile after 15 hours (times are morphological, the simulated time
times the acceleration factor).}
\end{figure}

\begin{figure}
\begin{center}
\includegraphics[width=0.9\textwidth]{suspended_load_test3.png}
\end{center}
\caption{Test 3: suspended load along the flume at the end of the run,
scaled by the inflow supply.}
\end{figure}

\IfFileExists{bed_profile_test1.png}{
\begin{figure}
\begin{center}
\includegraphics[width=0.9\textwidth]{bed_profile_test1.png}
\end{center}
\caption{Test 1, side slopes 1:3 (long run).}
\end{figure}
}{}

\IfFileExists{bed_profile_test2.png}{
\begin{figure}
\begin{center}
\includegraphics[width=0.9\textwidth]{bed_profile_test2.png}
\end{center}
\caption{Test 2, side slopes 1:7 (long run).}
\end{figure}
}{}

\endinput
